\vfill\pagebreak
\section{Naming the arguments}
\label{sec:functions:named_arguments}

In a call such as \token{printDate(4, 7, 1776)}, nothing says which number is
the day and which the month. The answer depends on the declaration of
\token{printDate}, and a reader used to the other convention reads the fourth of
July as the seventh of April. Two arguments of the same type can be swapped by
mistake, and the compiler cannot notice, since both orders have the right types.

\subsection{Named arguments}

An argument can name the parameter it is for, with the name of the parameter,
the arrow \token{->}, and the value: \token{printDate(day-> 4, month-> 7,
year-> 1776)}. A named argument goes to the parameter it names, wherever it is
written in the call, so the named arguments can come in any order. The arguments
without a name, called \textit{positional}, fill the remaining parameters in
order. Line~10 of the next listing names the year and the month, and the
positional \token{4} goes to the one parameter left, \token{day}.

\lstinputlisting[style=coloredverbatim, caption={Named arguments and a default value, \textit{examples/functions/dates.yr}}, label=lst:functions:dates]{examples/functions/dates.yr}
\vspace{-10pt}%
\begin{lstlisting}[style=bashVerb]
4/7/1776
4/7/1776
4/7/1776
14-7-1789
14.7.1789
\end{lstlisting}

The first three calls give the same arguments to the same parameters, and print
the same line. A parameter must get exactly one argument, so naming one twice,
or naming one that a positional argument has already filled, is refused.

\subsection{Default values}

A parameter can have a \textit{default value}, written after its type with
\token{=}, like the initial value of a variable. A call that gives no argument
for that parameter gets the default. \token{separator}, on line~3, defaults to
\token{'/'}: the calls of lines~8 to~10 do not mention it, and their dates are
printed with slashes.

A call that wants another value gives it \emph{by name}, as on lines~11 and~12.
A parameter with a default value is never filled by a positional argument, so the
last parameter of a declaration can have a default without making the others
optional, and a date is never read as a separator by accident.

\begin{lstlisting}[style=coloredverbatimError, escapechar=@, caption=A default value given by position]
use std::io;

fn printDate(day: i32, month: i32, year: i32, separator: c8 = '/') {
  println(day, separator, month, separator, year);
}

fn main() {
  printDate(14, 7, 1789, @\hb{'-'}@); // error
}
\end{lstlisting}

The compiler lists two reasons: \errmsg{E4212}{3 parameters are expected, but 4
  are provided}, since only the parameters without a default are filled by
position, and \errmsg{E4283}{the parameter separator has a default value, it can
  only be passed by name}.

\subsection{Arguments that must be named}

A named argument protects a call only if its author thinks of naming it. A
function can insist, by writing \token{@} in front of a parameter's name: that
parameter can then only be given by name. Declared with
\token{@day}, \token{@month} and \token{@year}, \token{printDate} can no longer
be called with three bare numbers, and every call says which date it means.

\begin{lstlisting}[style=coloredverbatimError, escapechar=|, caption=Parameters that must be named]
use std::io;

fn printDate(@day: i32, @month: i32, @year: i32) {
  println(day, "/", month, "/", year);
}

fn main() {
  printDate(day-> 14, month-> 7, year-> 1789); // ok
  printDate(|\hb{14, 7, 1789}|);                       // error
}
\end{lstlisting}

The second call is refused. The compiler counts
\errmsg{E4212}{0 parameters are expected, but 3 are provided}, since no parameter
can be filled by position, then gives one line per parameter:
\errmsg{E4284}{the parameter day can only be passed by name}, and the same for
\token{month} and \token{year}. The
\token{@} belongs to the declaration only; the body names the parameter
\token{day}, without it.

Whether a parameter should be named is a question about its calls. When the
function has a single parameter, or when the order of the arguments is obvious,
as in \token{rectangle(width, height)}, names only make the calls longer. When
several parameters have the same type and no natural order, as a day, a month
and a year, \token{@} rules out a whole class of mistakes. The calendar of
Section~\ref{sec:functions:calendar} uses it for exactly that.

\subsection{Reading, explained}
\label{sec:functions:read}

Chapter~\ref{chap:control} asked for \token{read!i32(ask-> "Your guess: ")} to
be taken as a fixed formula. Its second half is now ordinary. \token{read} is a
function of \token{std::io}, and its declaration is, in essence:

%% check: skip -- the real declaration needs templates; this is its shape
\begin{lstlisting}[style=coloredverbatim]
fn read(ask: [c8] = "")-> T
\end{lstlisting}

\token{ask} is a parameter with a default value, the empty string, so
\token{read!i32()} prints nothing before waiting, and a call that wants a prompt
gives it by name, \token{ask-> "Your guess: "}. The \token{T} is the part that
stays unexplained: \token{read} can return a value of many types, and
\token{!i32} is how the call chooses one. Functions that work on any type are
called \textit{templates}, and have a chapter of their own,
Chapter~\ref{chap:templates}.

\subsection{Calling with a dot}
\label{sec:functions:ucs}

A call can also be written with its first argument in front of the function
name, followed by a dot: \token{n.square()} is the call \token{square(n)}, and
\token{a.add(b)} is \token{add(a, b)}. The compiler rewrites the first form into
the second, so it works with any function, without changing its declaration.
This way of calling is the \textit{uniform call syntax}.

\begin{lstlisting}[style=coloredverbatim, caption=Calls written with a dot, label=lst:functions:ucs]
use std::io;

fn square(x: i32)-> i32 {
  x * x
}

fn add(a: i32, b: i32)-> i32 {
  a + b
}

fn main() {
  let n = 3;
  println(square(add(square(n), 1)));
  println(n.square().add(1).square());
}
\end{lstlisting}
\vspace{-10pt}%
\begin{lstlisting}[style=bashVerb]
100
100
\end{lstlisting}

Both lines compute the same thing: square 3, add 1, and square the result.
Nested calls, on line~13, are read from the inside out, starting with the
innermost \token{square(n)}. The chain on line~14 is read from left to right, in
the order the operations happen. The dot works after a literal as well, with one
catch: \token{3.square()} is read as the floating point number \token{3.} followed
by a name, and the literal must be put between parentheses,
\token{(3).square()}.
